\section{Density-FB}
\label{sec:method}

Density-FB models the normalized successor measure as a conditional density over future states. We first derive its categorical Bellman update, then use the learned successor distribution to construct the value signal for policy learning and downstream reward transfer.

\subsection{Learning normalized successor distributions}
\label{sec:method-density}

Since $M^\pi(s,a,\mathcal S)=1/(1-\gamma)$, we define the normalized successor measure as $\nu^\pi(s,a,\cdot):=(1-\gamma)M^\pi(s,a,\cdot)$. 
It is a probability measure over future states, and Equation~\eqref{eq:sm-bellman} gives the Bellman recursion
\begin{equation}
    \nu^\pi(s,a,\cdot) = (1-\gamma)P(\cdot\mid s,a) + \gamma\E_{\substack{s'\sim P(\cdot\mid s,a)\\a'\sim\pi(\cdot\mid s')}}\left[\nu^\pi(s',a',\cdot)\right].
    \label{eq:normalized-bellman}
\end{equation}
Thus, $1-\gamma$ of the successor probability comes from the immediate transition, while the remaining $\gamma$ is propagated from subsequent states.

For the policy indexed by $z$, let
\[
    h_z(s,a,s^+) := \frac{d\nu^{\pi_z}(s,a,\cdot)}{d\rho}(s^+)
\]
denote its normalized successor-density ratio. 
Whereas standard FB directly approximates the density of the unnormalized successor measure with a bilinear factorization, Density-FB uses the same factorization as a log-density score,
\begin{equation}
    f_\theta(s,a,z,s^+) := \frac{F_\theta(s,a,z)^\top B_\theta(s^+)}{\tau}, \qquad h_\theta(s,a,z,s^+) := \frac{\exp f_\theta(s,a,z,s^+)}{\E_{\tilde s^+\sim\rho}\left[\exp f_\theta(s,a,z,\tilde s^+)\right]}.
    \label{eq:density-model}
\end{equation}
Hence $F_\theta^\top B_\theta/\tau$ is a log-density score, while its normalized exponential models the successor-density ratio $h_z$.

We approximate the normalization in Equation~\eqref{eq:density-model} using the next states in a minibatch as samples from $\rho$. 
For $\{(s_i,a_i,s_i')\}_{i=1}^{n}$, define
\begin{equation}
    p_{ij} := \operatorname{softmax}_{j}\left(f_\theta(s_i,a_i,z_i,s_j')\right).
    \label{eq:density-softmax}
\end{equation}
Each row $p_i$ is a categorical distribution over the empirical candidate states $\{s_j'\}_{j=1}^{n}$.

To construct its Bellman target, let $a_i'=\sg[\pi_\omega(s_i',z_i)]$ and define the target-network prediction
\begin{equation}
    \bar p_{ij}' := \operatorname{softmax}_{j}\left(\frac{\bar F(s_i',a_i',z_i)^\top\bar B(s_j')}{\tau}\right).
    \label{eq:target-density-softmax}
\end{equation}
Equation~\eqref{eq:normalized-bellman} then yields the categorical Bellman target
\begin{equation}
    p_{ij}^{\mathrm{tgt}} := \sg\left[(1-\gamma)\mathbf 1[j=i]+\gamma\bar p_{ij}'\right].
    \label{eq:bellman-target}
\end{equation}
We fit the predicted distribution with cross-entropy,
\begin{equation}
    \mathcal L_{\mathrm{occ}} := -\frac{1}{n}\sum_{i=1}^{n}\sum_{j=1}^{n}p_{ij}^{\mathrm{tgt}}\log p_{ij}.
    \label{eq:occupancy-loss}
\end{equation}
This update parallels the categorical Bellman update in TD-InfoNCE~\citep{zheng-iclr24a}; here, we use it to learn the $z$-conditioned family of successor distributions required for FB-style control.

\subsection{Policy learning from successor distributions}
\label{sec:method-control}

A normalized successor density determines the value of any state-based reward by integration. 
For a policy $\pi_z$, define
\begin{equation}
    q_r^{\pi_z}(s,a) := \E_{s^+\sim\rho}\left[h_z(s,a,s^+)r(s^+)\right] = (1-\gamma)Q_r^{\pi_z}(s,a).
    \label{eq:normalized-value}
\end{equation}
Density-FB estimates this quantity by replacing the true density ratio with the learned model,
\begin{equation}
    q_{\theta,r}(s,a,z) := \E_{s^+\sim\rho}\left[h_\theta(s,a,z,s^+)r(s^+)\right].
    \label{eq:reward-readout-population}
\end{equation}
For a finite candidate pool $\{s_j^+\}_{j=1}^{P}$, we use
\begin{equation}
    q_{\theta,r}(s,a,z) \approx \sum_{j=1}^{P}p_j(s,a,z)r(s_j^+), \qquad p_j(s,a,z):=\operatorname{softmax}_{j}\left(f_\theta(s,a,z,s_j^+)\right).
    \label{eq:reward-readout}
\end{equation}
Thus, rather than reading value directly from $F(s,a,z)^\top z$ as in standard FB, Density-FB evaluates an action by integrating the requested reward under its predicted successor distribution.

During pretraining, each task index $z$ is paired with a reward $r_z$. 
In the goal-conditioned setting, for example, $z=g$ and $r_g$ rewards reaching the sampled goal. 
Let $\hat q_{s,z}:=q_{\theta,r_z}(s,\pi_\omega(s,z),z)$. We train a behavior-regularized actor with
\begin{equation}
    \mathcal L_{\pi} := -\frac{\E_{s,z}\left[\hat q_{s,z}\right]}{\sg\left[\E_{s,z}\left|\hat q_{s,z}-\E[\hat q_{s,z}]\right|+\varepsilon\right]}+\beta\mathcal L_{\mathrm{BC}}(\pi_\omega).
    \label{eq:actor-loss}
\end{equation}
The normalization controls the scale of the value term relative to behavior cloning. 
During the actor update, gradients are propagated through the action input to $F_\theta$.

\subsection{Zero-shot reward specification}
\label{sec:method-task-interface}

Standard FB obtains a downstream task vector from the linear identity
\begin{equation}
    \int_{\mathcal S} r(s^+)F(s,a,z)^\top B(s^+)\,d\rho(s^+) = F(s,a,z)^\top\E_\rho\left[r(s^+)B(s^+)\right],
    \label{eq:fb-linear-identity}
\end{equation}
which motivates the reward embedding $z_r\propto\E_\rho[rB]$. 
Density-FB does not preserve this identity because $F^\top B$ now parameterizes a log-density score rather than the successor density itself.

Moreover, the normalized model is invariant to additive shifts of the backward representation, whereas the standard FB reward moment is not:
\begin{equation}
    \operatorname{softmax}_{s^+}\left(\frac{F^\top(B(s^+)+b)}{\tau}\right)=\operatorname{softmax}_{s^+}\left(\frac{F^\top B(s^+)}{\tau}\right), \qquad \E_\rho[r(B+b)]=\E_\rho[rB]+b\,\E_\rho[r].
    \label{eq:task-gauge}
\end{equation}
Thus two backward representations that define the same Density-FB successor distribution can induce different values of $\E_\rho[rB]$; the standard FB reward embedding is therefore not an identifiable task interface for the normalized model. 
We provide the corresponding derivation in Appendix~\ref{app:fb-task-map}.

Instead, we keep $z$ in the task space used during pretraining and evaluate candidate policies directly under the downstream reward. 
For an unseen reward $r$, define
\begin{equation}
    J_{\theta,r}(z) := \E_{s\sim\mu_{\mathrm{sel}}}\left[q_{\theta,r}(s,\pi_\omega(s,z),z)\right], \qquad \hat z_r\in\arg\max_{z\in\mathcal Z_{\mathrm{cand}}}J_{\theta,r}(z).
    \label{eq:policy-search}
\end{equation}
Here $\mathcal Z_{\mathrm{cand}}$ contains task indices covered during pretraining and $\mu_{\mathrm{sel}}$ denotes the state distribution used for policy selection. 
During inference, all learned parameters remain frozen and we execute $\pi_\omega(\cdot,\hat z_r)$. 
This selects the highest-valued policy within the learned candidate family; 
unlike standard FB, it does not assume that an arbitrary reward can be mapped linearly to an optimal policy index.

\paragraph{Algorithm.}
Algorithm~\ref{alg:dfb} summarizes Density-FB training and zero-shot deployment.

\begin{algorithm}[t]
\caption{Density-FB}
\label{alg:dfb}
\begin{algorithmic}[1]
\Require Offline dataset $\mathcal D$, successor model $F_\theta,B_\theta$, target models $\bar F,\bar B$, actor $\pi_\omega$, discount $\gamma$, temperature $\tau$, BC coefficient $\beta$
\For{each training iteration}
    \State Sample transitions $\{(s_i,a_i,s_i')\}_{i=1}^{n}\sim\mathcal D$
    \State Sample task/reward pairs $\{(z_i,r_{z_i})\}_{i=1}^{n}$
    \State Compute successor probabilities $p_{ij}$ using Equation~\eqref{eq:density-softmax}
    \State $a_i'\gets\sg[\pi_\omega(s_i',z_i)]$
    \State Compute target probabilities $\bar p_{ij}'$ using Equation~\eqref{eq:target-density-softmax}
    \State $p_{ij}^{\mathrm{tgt}}\gets\sg[(1-\gamma)\mathbf 1[j=i]+\gamma\bar p_{ij}']$
    \State Update $F_\theta,B_\theta$ by minimizing $\mathcal L_{\mathrm{occ}}$ in Equation~\eqref{eq:occupancy-loss}
    \State Estimate $q_{\theta,r_{z_i}}(s_i,\pi_\omega(s_i,z_i),z_i)$ using Equation~\eqref{eq:reward-readout}
    \State Update $\pi_\omega$ by minimizing $\mathcal L_\pi$ in Equation~\eqref{eq:actor-loss}
    \State Update $\bar F,\bar B$ by exponential moving average
\EndFor
\Statex
\Statex \textbf{Zero-shot deployment}
\State Given downstream reward $r$, select $\hat z_r$ from $\mathcal Z_{\mathrm{cand}}$ using Equation~\eqref{eq:policy-search}
\State Execute the frozen policy $\pi_\omega(\cdot,\hat z_r)$
\end{algorithmic}
\end{algorithm}